% !TEX program = xelatex
\documentclass[11pt]{article}

\usepackage{fontspec}
\setmainfont{CMU Serif}

\usepackage[margin=1in]{geometry}
\usepackage{amsmath,amssymb}
\usepackage{mathtools}
\usepackage{physics}
\usepackage{hyperref}
\usepackage{microtype}
\usepackage{enumitem}

\setlength{\parindent}{0pt}
\setlength{\parskip}{0.7\baselineskip}

\title{README: Многоуровневая архитектура GRA Мета-обнулёнка в мультиверсе}
\author{}
\date{}

\begin{document}
\maketitle

\section{Назначение документа}

Этот файл фиксирует \emph{формальную} (математическую) версию концепции ``GRA Мета-обнулёнка'', расширенной на многоуровневую (иерархическую) структуру \emph{мультиверса}: множество мета-систем, согласуемых рекурсивно от локальных доменов до глобального уровня.

Цель README --- дать компактные определения, единообразные обозначения, ключевые функционалы, условие ``полного обнуления'' (нулевой пены на всех уровнях) и основные выводы.

\section{Иерархия уровней и мультииндексация}

\subsection{Уровни мультиверса}

Пусть задана глубина иерархии $K\in\mathbb{N}$.

\begin{itemize}[leftmargin=*, itemsep=2pt]
  \item \textbf{Уровень $0$ (локальный)}: обнуление внутри отдельных доменов.
  \item \textbf{Уровень $1$ (мета)}: согласование доменов в пределах одной мета-системы.
  \item \textbf{Уровень $2$ (мета-мета)}: согласование различных мета-систем.
  \item \dots
  \item \textbf{Уровень $K$ (глобальный)}: согласование всей иерархии.
\end{itemize}

\subsection{Мультииндекс}

Каждому объекту (подсистеме/домену/мета-системе) сопоставляется мультииндекс
\[
\mathbf{a}=(a_0,a_1,\dots,a_k), \qquad 0\le k\le K,
\]
где $a_0$ индексирует домен, $a_1$ --- мета-систему и т.д. Размерность (``уровень'') мультииндекса обозначим
\[
\dim(\mathbf{a})\coloneqq k.
\]

Набор всех индексов уровня $l$ обозначим $\mathcal{I}_l\coloneqq\{\mathbf{a}:\dim(\mathbf{a})=l\}$, а полный набор всех индексов --- $\mathcal{I}\coloneqq\bigcup_{l=0}^K\mathcal{I}_l$.

\subsection{Отношение вложенности подсистем}

Пишем $\mathbf{b}\prec\mathbf{a}$, если $\mathbf{b}$ является непосредственной или транзитивной подсистемой $\mathbf{a}$ (например, ``дочерний домен'' внутри мета-системы). Формально, это частичный порядок на $\mathcal{I}$, согласованный с уровнем: если $\mathbf{b}\prec\mathbf{a}$, то $\dim(\mathbf{b})<\dim(\mathbf{a})$.

\section{Пространства состояний}

\subsection{Локальные пространства}

Каждому мультииндексу $\mathbf{a}\in\mathcal{I}$ сопоставляется гильбертово пространство состояний $\mathcal{H}^{(\mathbf{a})}$ и вектор состояния $\ket{\Psi^{(\mathbf{a})}}\in\mathcal{H}^{(\mathbf{a})}$.

\subsection{Тензорная агрегация по уровням}

Для уровня $l$ вводим агрегированное пространство:
\[
\mathcal{H}^{(l)} \coloneqq \bigotimes_{\mathbf{a}\in\mathcal{I}_l} \mathcal{H}^{(\mathbf{a})}.
\]

Полное пространство мультиверса:
\[
\mathcal{H}_{\text{multiverse}} \coloneqq \bigotimes_{l=0}^{K} \mathcal{H}^{(l)}.
\]

Состояние всей иерархии будем записывать как семейство
\[
\mathbf{\Psi} \coloneqq \{\Psi^{(\mathbf{a})}\}_{\mathbf{a}\in\mathcal{I}}.
\]

\section{Иерархия целей и проекторов}

\subsection{Цели уровня}

Каждому уровню соответствует своя цель согласования:
\begin{itemize}[leftmargin=*, itemsep=2pt]
  \item локальные цели $G_0^{(\mathbf{a})}$, $\mathbf{a}\in\mathcal{I}_0$;
  \item мета-цели $G_1^{(\mathbf{a})}$, $\mathbf{a}\in\mathcal{I}_1$;
  \item \dots
  \item глобальная цель $G_K$ (или $G_K^{(\mathbf{a})}$ при необходимости индексировать верхний уровень).
\end{itemize}

\subsection{Проекторы на ``пространства решений''}

Цели кодируются проекторами $\mathcal{P}_{G_l^{(\mathbf{a})}}$ (ортопроекторами) на подпространство состояний, удовлетворяющих цели уровня $l$ в подсистеме $\mathbf{a}$. Для краткости, когда индекс подсистемы не важен, будем писать $\mathcal{P}_{G_l}$.

\section{``Пена'' как мера рассогласования}

\subsection{Интуиция}

Пена $\Phi$ измеряет степень \emph{взаимного несовпадения} (неортогональности/интерференции вне согласованного базиса) между компонентами состояния различных подсистем одного уровня относительно проектора на решение общей цели.

\subsection{Определение пены уровня $l$}

Для агрегированного состояния уровня $l$ и цели $G_l$:
\[
\Phi^{(l)}\bigl(\Psi^{(l)},G_l\bigr)
\;\coloneqq\;
\sum_{\substack{\mathbf{a},\mathbf{b}\in\mathcal{I}_l\\ \mathbf{a}\neq\mathbf{b}}}
\abs{\mel{\Psi^{(\mathbf{a})}}{\mathcal{P}_{G_l}}{\Psi^{(\mathbf{b})}}}^2.
\]

Свойства:
\begin{itemize}[leftmargin=*, itemsep=2pt]
  \item $\Phi^{(l)}\ge 0$ (сумма квадратов модулей).
  \item $\Phi^{(l)}=0$ тогда и только тогда, когда все ``внедиагональные'' элементы между различными подсистемами уровня $l$ обнулены в базисе, согласованном с $\mathcal{P}_{G_l}$.
\end{itemize}

\section{Супер-мета-функционал (глобальный критерий)}

\subsection{Весовая агрегация уровней}

Глобальный функционал мультиверса определим как сумму по уровням с весами $\Lambda_l$:
\[
J_{\text{multiverse}}(\mathbf{\Psi})
\;\coloneqq\;
\sum_{l=0}^{K}\Lambda_l\sum_{\mathbf{a}\in\mathcal{I}_l} J^{(l)}\bigl(\Psi^{(\mathbf{a})}\bigr).
\]

Удобная (и устойчиво затухающая) параметризация весов:
\[
\Lambda_l\coloneqq \lambda_0\alpha^l,\qquad \lambda_0>0,\; 0<\alpha<1.
\]

\subsection{Рекурсивная структура уровневого функционала}

Базовый уровень:
\[
J^{(0)}\bigl(\Psi^{(\mathbf{a})}\bigr) \coloneqq J_{\text{loc}}\bigl(\Psi^{(\mathbf{a})};\,G_0^{(\mathbf{a})}\bigr).
\]

Рекурсия для $l\ge 1$:
\[
J^{(l)}\bigl(\Psi^{(\mathbf{a})}\bigr)
\;\coloneqq\;
\sum_{\substack{\mathbf{b}\prec\mathbf{a}\\ \dim(\mathbf{b})=l-1}}
J^{(l-1)}\bigl(\Psi^{(\mathbf{b})}\bigr)
\;+
\Phi^{(l)}\bigl(\Psi^{(\mathbf{a})},\,G_l^{(\mathbf{a})}\bigr).
\]

Смысл: функционал верхнего уровня складывает ``стоимость'' всех нижележащих подсистем и добавляет штраф за рассогласование текущего уровня.

\section{Условия полного мультиверсного обнуления}

\subsection{Условия}

Рассмотрим три условия, которые делают возможным \emph{глобально согласованное} состояние с нулевой пеной на каждом уровне.

\begin{enumerate}[leftmargin=*, itemsep=4pt]
  \item \textbf{Полная коммутативность проекторов:}
  \[
    \comm{\mathcal{P}_{G_l^{(\mathbf{a})}}}{\mathcal{P}_{G_m^{(\mathbf{b})}}}=0
    \quad \forall\,\mathbf{a},\mathbf{b},\; \forall\,l,m.
  \]

  \item \textbf{Согласованность иерархии (факторизация цели):}
  \[
    \mathcal{P}_{G_l^{(\mathbf{a})}}
    = \bigotimes_{\mathbf{b}\prec\mathbf{a}} \mathcal{P}_{G_{l-1}^{(\mathbf{b})}}
    \quad \forall\, l\ge 1.
  \]

  \item \textbf{Полнота пространства (достаточная размерность):}
  \[
    \dim\bigl(\mathcal{H}_{\text{multiverse}}\bigr)\ge \prod_{l=0}^{K} N_l,
  \]
  где $N_l$ --- число подсистем на уровне $l$.
\end{enumerate}

\subsection{Теорема (мультиверсное обнуление)}

\textbf{Теорема.}
Если выполнены условия (1)--(3), то существует состояние $\mathbf{\Psi}^*$ такое, что
\[
\Phi^{(l)}\bigl(\Psi^{(l)*},G_l\bigr)=0,\qquad \forall\, l\in\{0,1,\dots,K\}.
\]

\subsection{Схема доказательства}

\textbf{База ($l=0$):} из локальной теоремы обнуления следует существование $\Psi^{(0)*}$, согласованного с $\mathcal{P}_{G_0}$.

\textbf{Переход ($l-1\to l$):}
если все подсистемы уровня $l-1$ выбраны как собственные векторы соответствующих проекторов, то по согласованности (2) тензорное произведение этих собственных векторов является собственным вектором проектора уровня $l$.
Коммутативность (1) гарантирует общий собственный базис для всей иерархии проекторов; в таком базисе внедиагональные элементы между различными подсистемами исчезают, т.е. $\Phi^{(l)}=0$.

\section{Алгоритмическая реализация (рекурсия + мета-согласование)}

\subsection{Рекурсивный шаг}

Опишем абстрактную процедуру ``обнуления уровня'':
\begin{enumerate}[leftmargin=*, itemsep=4pt]
  \item если $l=0$, применить локальный алгоритм оптимизации для $J_{\text{loc}}$;
  \item иначе разложить состояние уровня $l$ на подсистемы уровня $l-1$;
  \item рекурсивно обнулить каждую подсистему;
  \item собрать состояние обратно;
  \item выполнить оптимизацию/градиентный спуск по пене $\Phi^{(l)}$ до порога $\varepsilon_l$.
\end{enumerate}

\subsection{Параллельные обновления}

Для отдельного состояния $\Psi^{(\mathbf{a})}$ уровня $l$ формально полезно учитывать как вклад собственного уровня, так и вклад ``родительских'' уровней:
\[
\pdv{J_{\text{multiverse}}}{\Psi^{(\mathbf{a})}}
= \Lambda_l\,\pdv{\Phi^{(l)}}{\Psi^{(\mathbf{a})}}
+ \sum_{\mathbf{b}\succ\mathbf{a}} \Lambda_{l+1}\,\pdv{\Phi^{(l+1)}}{\Psi^{(\mathbf{a})}}.
\]

Отсюда получается схема параллельного шага:
\[
\Psi^{(\mathbf{a})}(t+1)=\Psi^{(\mathbf{a})}(t)-\eta\left(
\Lambda_l\,\pdv{\Phi^{(l)}}{\Psi^{(\mathbf{a})}}
+ \sum_{\mathbf{b}\succ\mathbf{a}} \Lambda_{l+1}\,\pdv{\Phi^{(l+1)}}{\Psi^{(\mathbf{a})}}
\right).
\]

\section{Следствия и выводы}

\subsection{Структура решения при полном обнулении}

В состоянии полного обнуления верхнего уровня:
\[
\mathcal{P}_{G_K}\ket{\Psi_{\text{multiverse}}^*}=\ket{\Psi_{\text{multiverse}}^*}.
\]

И для любого уровня $l<K$ и любой подсистемы $\mathbf{a}\in\mathcal{I}_l$:
\[
\mathcal{P}_{G_l^{(\mathbf{a})}}\ket{\Psi^{(\mathbf{a})*}}=\ket{\Psi^{(\mathbf{a})*}}.
\]

\subsection{Единственность (с точностью до симметрий)}

Если проекторы задают общую систему ограничений, то решение фиксируется с точностью до:
\begin{itemize}[leftmargin=*, itemsep=2pt]
  \item глобальных фаз (на каждом уровне/в каждой компоненте);
  \item унитарных преобразований, коммутирующих со всей иерархией проекторов.
\end{itemize}

Интерпретационно это описывается как ``абсолютный когнитивный вакуум'': отсутствие артефактов интерпретации на всех уровнях.

\subsection{Оценка сложности}

При предположении $N$ подсистем на каждом уровне и геометрическом затухании весов $\Lambda_l\propto\alpha^l$:
\[
\text{Complexity}
= O\left(\sum_{l=0}^{K} N^2\alpha^l\right)
= O\left(N^2\,\frac{1-\alpha^{K+1}}{1-\alpha}\right).
\]

В пределе $K\to\infty$ при $0<\alpha<1$:
\[
\text{Complexity}=O\left(\frac{N^2}{1-\alpha}\right),
\]
то есть сложность остаётся полиномиальной при произвольной глубине иерархии (при фиксированных $N$ и $\alpha$).

\subsection{Предельная формула (бесконечная иерархия)}

Если существует предел последовательности решений $\{\Psi_K^*\}_{K\in\mathbb{N}}$, то можно определить предельное состояние:
\[
\Psi_{\infty}^*\coloneqq \lim_{K\to\infty}\Psi_K^*.
\]

Тогда для любого фиксированного $l\in\mathbb{N}$:
\[
\Phi^{(l)}\bigl(\Psi_{\infty}^*,G_l\bigr)=0.
\]

Итоговая формула:
\[
\boxed{
\lim_{K\to\infty}\Psi_K^* = \Psi_{\infty}^*
\quad\text{такое, что}\quad
\bigcap_{l=0}^{\infty}\ker\bigl(\Phi^{(l)}\bigr)=\{\Psi_{\infty}^*\}
}.
\]

\subsection{Практическая памятка (что настраивать)}

\begin{itemize}[leftmargin=*, itemsep=2pt]
  \item Выберите глубину $K$ (сколько уровней абстракции действительно нужно моделировать).
  \item Задайте семейство целей $\{G_l^{(\mathbf{a})}\}$ и проверьте совместимость/коммутативность проекторов.
  \item Настройте веса $\Lambda_l$ (обычно геометрическое затухание) и пороги $\varepsilon_l$.
  \item Реализуйте локальный решатель $J_{\text{loc}}$ и мета-оптимизацию пены на каждом уровне.
\end{itemize}

\end{document}